% B4 插补消融的三方案对照文本，供 D 在结果章落位（表注或正文句）。
% 数字与 results/b4_missing_comparison.csv 一致（cv10, ddof=0）。
% 说明：数据章 sec3b 已含分点决策与表格（tab:b_imputation），此处为结果章备选写法。

% ── 方案一（推荐，对应批注 #12/#13）：三方案对照句，与 A6 同款书写 ──
% As a controlled comparison behind the imputation choice on sick, written
% in the same style as the encoding ablation on credit-g: fold-internal
% imputation scored F1 $0.830\pm0.049$ with AUC $0.974$ (adopted);
% dropping incomplete rows $0.843\pm0.069$ with $0.964$ --- within one
% standard deviation, at the cost of 1{,}129 rows including 19 of the 231
% positives; treating ``missing'' as its own category $0.828$. Decision:
% impute, inside the pipeline (cv10, ddof$=0$;
% results/b4\_missing\_comparison.csv).

% ── 方案二：表脚注（若维持原 D5 方案）──
% \footnote{The imputation decision was validated as a controlled comparison
% on sick: fold-internal imputation scored F1 $0.830\pm0.049$ with AUC
% $0.974$, versus $0.843\pm0.069$ with AUC $0.964$ for dropping incomplete
% rows --- a difference within one standard deviation that sacrifices
% 1{,}129 rows (including 19 of the 231 positives); treating ``missing''
% as a category scored $0.828$. Imputation was therefore locked in the
% preprocessing spec (cv10, ddof$=0$; results/b4\_missing\_comparison.csv).}

% 提示：§3.2.2 v1.1 已按 D5+#11 收紧（docs/B8_sec3b_sick_draft_v1.0.tex）；
% D 落位任一方案后，D5 与批注 #12/#13 同时闭环。
